\chapter{Complexity --- Three SAT To Clique}
%%% generated by the blueprint pipeline from TCSlib.Complexity.NPReductions.ThreeSATToClique
\section{Overview}

This module formalises the classical NP-hardness reduction from 3-SAT to the Clique
problem.  Given a 3-CNF formula $f$ with $m$ clauses, a \emph{conflict graph} is
constructed whose vertices are pairs (clause index, literal position) and whose edges
connect vertices from different clauses whose literals are not complementary; the main
results establish that $f$ is satisfiable if and only if the conflict graph contains
an $m$-clique.

%%%%%%%%%%%%%%%%%%%%%%%%%%%%%%%%%%%%%%%%%%%%%%%%%%%%%%%%%%%%%%%%%%%%%%%%%%%%%%%
\section{Declarations}

\begin{definition}[Clause]
\lean{ThreeSATToClique.Clause}
\label{ThreeSATToClique.Clause}
\leanok
A clause is a list of literals over a variable type $V$; it represents a
disjunction of those literals.
\end{definition}

\begin{definition}[CNF formula]
\lean{ThreeSATToClique.CNFFormula}
\label{ThreeSATToClique.CNFFormula}
\leanok
\uses{ThreeSATToClique.Clause}
A CNF formula over $V$ is a list of clauses, representing the conjunction of those
clauses.
\end{definition}

\begin{definition}[Literal evaluation]
\lean{ThreeSATToClique.evalLiteral}
\label{ThreeSATToClique.evalLiteral}
\leanok
Given a truth assignment $\alpha : V \to \mathrm{Prop}$, the evaluation of a
literal $\ell$ under $\alpha$ is $\alpha(v)$ when $\ell = \texttt{pos}\,v$ and
$\neg\,\alpha(v)$ when $\ell = \texttt{neg}\,v$.
\end{definition}

\begin{definition}[Clause satisfaction]
\lean{ThreeSATToClique.clauseSatisfied}
\label{ThreeSATToClique.clauseSatisfied}
\leanok
\uses{ThreeSATToClique.Clause, ThreeSATToClique.evalLiteral}
A clause $c$ is \emph{satisfied} by assignment $\alpha$ if at least one literal
$\ell \in c$ satisfies $\mathtt{evalLiteral}\,\alpha\,\ell$.
\end{definition}

\begin{definition}[CNF formula satisfaction]
\lean{ThreeSATToClique.formulaSatisfied}
\label{ThreeSATToClique.formulaSatisfied}
\leanok
\uses{ThreeSATToClique.CNFFormula, ThreeSATToClique.clauseSatisfied}
A CNF formula $f$ is \emph{satisfied} by $\alpha$ if every clause $c \in f$ is
satisfied by $\alpha$.
\end{definition}

\begin{definition}[CNF satisfiability]
\lean{ThreeSATToClique.isSatisfiable}
\label{ThreeSATToClique.isSatisfiable}
\leanok
\uses{ThreeSATToClique.CNFFormula, ThreeSATToClique.formulaSatisfied}
A CNF formula $f$ is \emph{satisfiable} if there exists an assignment $\alpha$
that satisfies $f$.
\end{definition}

\begin{definition}[3-clause]
\lean{ThreeSATToClique.Clause3}
\label{ThreeSATToClique.Clause3}
\leanok
A 3-clause over $V$ is a structure with exactly three literals $l_1, l_2, l_3 :
\mathtt{Literal}\,V$, representing a disjunction of precisely three literals.
\end{definition}

\begin{definition}[3-CNF formula]
\lean{ThreeSATToClique.Formula3}
\label{ThreeSATToClique.Formula3}
\leanok
\uses{ThreeSATToClique.Clause3}
A 3-CNF formula over $V$ is a list of 3-clauses, representing their conjunction.
\end{definition}

\begin{definition}[3-clause satisfaction]
\lean{ThreeSATToClique.clause3Satisfied}
\label{ThreeSATToClique.clause3Satisfied}
\leanok
\uses{ThreeSATToClique.Clause3, ThreeSATToClique.evalLiteral}
A 3-clause $c$ is satisfied by $\alpha$ if at least one of its three literals
$c.l_1$, $c.l_2$, $c.l_3$ evaluates to true under $\alpha$.
\end{definition}

\begin{definition}[3-CNF formula satisfaction]
\lean{ThreeSATToClique.formula3Satisfied}
\label{ThreeSATToClique.formula3Satisfied}
\leanok
\uses{ThreeSATToClique.Formula3, ThreeSATToClique.clause3Satisfied}
A 3-CNF formula $f$ is satisfied by $\alpha$ if every 3-clause $c \in f$ is
satisfied by $\alpha$.
\end{definition}

\begin{definition}[3-CNF satisfiability]
\lean{ThreeSATToClique.is3Satisfiable}
\label{ThreeSATToClique.is3Satisfiable}
\leanok
\uses{ThreeSATToClique.Formula3, ThreeSATToClique.formula3Satisfied}
A 3-CNF formula $f$ is satisfiable if there exists an assignment $\alpha$ such
that $\alpha$ satisfies $f$.
\end{definition}

\begin{definition}[Clique vertex]
\lean{ThreeSATToClique.CliqueVertex}
\label{ThreeSATToClique.CliqueVertex}
\leanok
A vertex of the conflict graph for a formula with $m$ clauses is a pair
$\langle c\_idx, l\_idx \rangle$ where $c\_idx : \mathrm{Fin}\,m$ is a clause index
and $l\_idx : \mathrm{Fin}\,3$ is a literal position within that clause.
\end{definition}

\begin{definition}[Literal extraction from a 3-clause]
\lean{ThreeSATToClique.getLitInClause}
\label{ThreeSATToClique.getLitInClause}
\leanok
\uses{ThreeSATToClique.Clause3}
Given a 3-clause $c$ and a position $p : \mathrm{Fin}\,3$, returns the $p$-th
literal of $c$: $c.l_1$ at position $0$, $c.l_2$ at position $1$, and $c.l_3$ at
position $2$.
\end{definition}

\begin{definition}[Literal named by a vertex]
\lean{ThreeSATToClique.getLitAt}
\label{ThreeSATToClique.getLitAt}
\leanok
\uses{ThreeSATToClique.CliqueVertex, ThreeSATToClique.Formula3, ThreeSATToClique.getLitInClause}
Given a 3-CNF formula $f$ and a vertex $v : \mathtt{CliqueVertex}\,f.\mathtt{length}$,
returns the literal at position $v.l\_idx$ within clause $f[v.c\_idx]$.  The type of
$v$ guarantees the index is always in bounds.
\end{definition}

\begin{definition}[Literal conflict]
\lean{ThreeSATToClique.literalsConflict}
\label{ThreeSATToClique.literalsConflict}
\leanok
Two literals $l_1$ and $l_2$ \emph{conflict} if one is the positive and the other
the negative occurrence of the same variable, i.e.\ $l_1 = \texttt{pos}\,v$ and
$l_2 = \texttt{neg}\,v$ (or vice versa) for some $v$.
\end{definition}

\begin{theorem}[Symmetry of literal conflict]
\lean{ThreeSATToClique.literalsConflict_symm}
\label{ThreeSATToClique.literalsConflict_symm}
\leanok
\uses{ThreeSATToClique.literalsConflict}
\difficulty{3}
Call two literals over a variable set *conflicting* when one is the positive and the
other the negative occurrence of one and the same variable. For any two literals $l_1$
and $l_2$, $l_1$ and $l_2$ conflict if and only if $l_2$ and $l_1$ conflict; that is,
the conflict relation is symmetric.
\end{theorem}

\begin{definition}[Conflict graph]
\lean{ThreeSATToClique.toCliqueGraph}
\label{ThreeSATToClique.toCliqueGraph}
\leanok
\uses{ThreeSATToClique.CliqueVertex, ThreeSATToClique.Formula3, ThreeSATToClique.getLitAt, ThreeSATToClique.literalsConflict, ThreeSATToClique.literalsConflict_symm}
The conflict graph of a 3-CNF formula $f$ is the simple graph on vertex set
$\mathtt{CliqueVertex}\,f.\mathtt{length}$ in which two vertices $u$ and $v$ are
adjacent if and only if they come from \emph{different} clauses ($u.c\_idx \ne
v.c\_idx$) and their respective literals do not conflict.
\end{definition}

\begin{definition}[$k$-clique existence]
\lean{ThreeSATToClique.hasClique}
\label{ThreeSATToClique.hasClique}
\leanok
A graph $G$ on vertex type $V$ \emph{has a $k$-clique} if there exists a finite set
$s \subseteq V$ of $k$ vertices that are pairwise adjacent in $G$.
\end{definition}

\begin{lemma}[True literals never conflict]
\lean{ThreeSATToClique.no_conflict_of_true}
\label{ThreeSATToClique.no_conflict_of_true}
\leanok
\uses{ThreeSATToClique.evalLiteral, ThreeSATToClique.literalsConflict}
\difficulty{3}
Let $\alpha$ be a truth assignment on a set $V$ of variables, and let $\ell_1,\ell_2$ be
literals over $V$, each of the form $\mathrm{pos}\,v$ (the variable $v$) or
$\mathrm{neg}\,v$ (its negation). If both $\ell_1$ and $\ell_2$ evaluate to true under
$\alpha$, then they do not conflict; that is, they are not the positive and negative
occurrences of one and the same variable.
\end{lemma}

\begin{lemma}[One clique vertex per clause]
\lean{ThreeSATToClique.clique_vertices_choose_one_per_clause}
\label{ThreeSATToClique.clique_vertices_choose_one_per_clause}
\leanok
\uses{ThreeSATToClique.CliqueVertex, ThreeSATToClique.Formula3, ThreeSATToClique.toCliqueGraph}
\difficulty{4}
Let $f$ be a 3-CNF formula with $m$ clauses, and let $s$ be a set of vertices of the
conflict graph of $f$. Suppose that $s$ has exactly $m$ elements and forms a clique,
meaning any two distinct vertices of $s$ are adjacent. Then for every clause index $i
\in \mathrm{Fin}\,m$ there is exactly one vertex $u \in s$ whose clause component equals
$i$.
\end{lemma}

\begin{lemma}[Vertex literal lies among its clause's three literals]
\lean{ThreeSATToClique.getLitAt_mem_clause}
\label{ThreeSATToClique.getLitAt_mem_clause}
\leanok
\uses{ThreeSATToClique.CliqueVertex, ThreeSATToClique.Formula3, ThreeSATToClique.getLitAt, ThreeSATToClique.getLitInClause}
\difficulty{3}
Let $f$ be a 3-CNF formula, let $v$ be a vertex of the conflict graph of $f$, and let
$i$ be a clause index of $f$ such that the clause component of $v$ equals $i$. Then the
literal named by $v$ is one of the three literals $l_1, l_2, l_3$ of the clause $f[i]$.
\end{lemma}

\begin{theorem}[Completeness of the 3-SAT to clique reduction]
\lean{ThreeSATToClique.ThreeSAT_to_Clique_completeness}
\label{ThreeSATToClique.ThreeSAT_to_Clique_completeness}
\leanok
\uses{ThreeSATToClique.CliqueVertex, ThreeSATToClique.Formula3, ThreeSATToClique.evalLiteral, ThreeSATToClique.getLitAt, ThreeSATToClique.hasClique, ThreeSATToClique.is3Satisfiable, ThreeSATToClique.no_conflict_of_true, ThreeSATToClique.toCliqueGraph}
\difficulty{5}
Let $f$ be a 3-CNF formula consisting of $m$ clauses, and let its conflict graph be the
simple graph on the vertices $\langle c, \ell\rangle$ (clause index $c$, literal
position $\ell$) in which two vertices are adjacent exactly when they belong to
different clauses and their literals do not conflict. If $f$ is satisfiable, then this
conflict graph has an $m$-clique.
\end{theorem}

\begin{theorem}[Soundness of the 3-SAT to clique reduction]
\lean{ThreeSATToClique.ThreeSAT_to_Clique_soundness}
\label{ThreeSATToClique.ThreeSAT_to_Clique_soundness}
\leanok
\uses{ThreeSATToClique.Formula3, ThreeSATToClique.clique_vertices_choose_one_per_clause, ThreeSATToClique.evalLiteral, ThreeSATToClique.getLitAt, ThreeSATToClique.getLitAt_mem_clause, ThreeSATToClique.hasClique, ThreeSATToClique.is3Satisfiable, ThreeSATToClique.toCliqueGraph}
\difficulty{5}
Let $f$ be a 3-CNF formula consisting of $m$ clauses. If the conflict graph of $f$ has
an $m$-clique, then $f$ is satisfiable.
\end{theorem}

\begin{theorem}[Correctness of the 3-SAT to Clique reduction]
\lean{ThreeSATToClique.ThreeSAT_to_Clique_equivalence}
\label{ThreeSATToClique.ThreeSAT_to_Clique_equivalence}
\leanok
\uses{ThreeSATToClique.Formula3, ThreeSATToClique.ThreeSAT_to_Clique_completeness, ThreeSATToClique.ThreeSAT_to_Clique_soundness, ThreeSATToClique.hasClique, ThreeSATToClique.is3Satisfiable, ThreeSATToClique.toCliqueGraph}
\difficulty{1}
Let $f$ be a 3-CNF formula with $m$ clauses. Then $f$ is satisfiable if and only if its
conflict graph has a clique of size $m$; that is, the graph whose vertices are the
literal positions of $f$, with two positions adjacent exactly when they lie in different
clauses and their literals do not conflict, contains $m$ pairwise-adjacent vertices
precisely when some truth assignment satisfies every clause of $f$.
\end{theorem}
